%% SYNTHETIC TEST UNIT generated by writing/tools/make_outline_book.py from FINAL_OUTLINE_v2 -- not book text
\chapter{引言：部分子图像、因子化的含义与严格性分级}\label{ch:intro}
本单元为测试用的合成文本（由 FINAL\_OUTLINE\_v2 自动生成），陈述内容从略。

\section{第 1 节（测试）}\label{sec:intro-parton}

\begin{definition}[硬过程与幂次展开]\label{def:intro-hard-process}
（测试用合成陈述：内容从略。）
% uses: none
\end{definition}

\begin{definition}[共线因子化陈述]\label{def:intro-factorization}
（测试用合成陈述：内容从略。）
\uses{def:intro-hard-process}
\end{definition}

\section{第 2 节（测试）}\label{sec:intro-statement}

\begin{remark}[部分子模型的非相干性论证（启发式）]\label{rem:intro-incoherence}
（测试用合成陈述：内容从略。）
% uses: none
\end{remark}

\begin{example}[树图手提包与 Callan–Gross 关系]\label{ex:intro-handbag}
（测试用合成陈述：内容从略。）
\uses{def:intro-factorization}
\end{example}

\section{第 3 节（测试）}\label{sec:intro-proof}

\begin{remark}["每一阶、领头幂次"的精确含义与幂次的锐利性]\label{rem:intro-sharpness}
（测试用合成陈述：内容从略。）
\uses{def:intro-factorization}
\end{remark}

\begin{definition}[严格性标签、限定语与等级传递]\label{def:intro-tags}
（测试用合成陈述：内容从略。）
非微扰地位：\cref{thm:dy-factorization}、\cref{asm:appN-collins}、\cref{asm:appO-mixing}、\cref{asm:intro}、\cref{asm:appN-brst-nonlocal}、\cref{lem:method-ward-multi}、\cref{thm:method-template}、\cref{thm:appN-structure}。
% uses: none
\end{definition}

\section{第 4 节（测试）}\label{sec:intro-tags}

\begin{definition}[误差类型]\label{def:intro-error-types}
（测试用合成陈述：内容从略。）
\uses{def:intro-tags}
\end{definition}

\begin{remark}[与文献中另一种分级的对照]\label{rem:intro-grade-dictionary}
（测试用合成陈述：内容从略。）
\uses{def:intro-tags}
\end{remark}

\section{第 5 节（测试）}\label{sec:intro-kinds}

\begin{assumption}[QCD 的 Wightman 结构 \tagC]\label{asm:intro-wightman}
（测试用合成陈述：内容从略。）
非微扰地位：\cref{asm:intro-composite}、\cref{asm:intro-np-ope}、\cref{asm:intro-growth}、\cref{prop:prelim-optical}、\cref{prop:ee-R-ImPi}、\cref{prop:ee-analytic}、\cref{prop:ee-smearing}、\cref{prop:dis-W-structure}、\cref{prop:dis-ope-rg}、\cref{prop:dis-dispersion-np}、\cref{thm:dis-fixed-n}、\cref{thm:dis-moments-np}、\cref{cor:dis-weak-xspace}、\cref{thm:dis-moments}、\cref{prop:dis-twist4}、\cref{prop:dis-tmc}、\cref{thm:dis-bjorken}、\cref{prop:pdf-invariance}、\cref{prop:pdf-support-hadron}、\cref{thm:pdf-moments-local}、\cref{thm:pdf-dis-xspace}、\cref{prop:pdf-sum-rules}、\cref{prop:pdf-timelike-dglap}、\cref{prop:coll-photon-sf}、\cref{prop:coll-luxqed}、\cref{thm:coll-g1-moments}、\cref{prop:coll-transversity-evolution}、\cref{thm:tmd-sivers-sign}、\cref{prop:excl-erbl}、\cref{prop:excl-polynomiality}、\cref{thm:heavy-inclusive-b}、\cref{thm:lamet-exp-decay}、\cref{lem:appM-forward-tube}、\cref{lem:appM-positivity}、\cref{def:intro-MA}。
地位：测试状态句。
文献：测试。
被使用于：\cref{asm:intro-composite}、\cref{asm:intro-np-ope}、\cref{asm:intro-growth}、\cref{def:intro-MA}、\cref{prop:prelim-optical}、\cref{prop:ee-R-ImPi}、\cref{prop:ee-analytic}、\cref{prop:ee-smearing}、\cref{prop:dis-W-structure}、\cref{prop:dis-ope-rg}、\cref{prop:dis-dispersion-np}、\cref{thm:dis-fixed-n}、\cref{thm:dis-moments-np}、\cref{cor:dis-weak-xspace}、\cref{thm:dis-moments}、\cref{prop:dis-twist4}、\cref{prop:dis-tmc}、\cref{thm:dis-bjorken}、\cref{prop:pdf-invariance}、\cref{prop:pdf-support-hadron}、\cref{thm:pdf-moments-local}、\cref{thm:pdf-dis-xspace}、\cref{prop:pdf-sum-rules}、\cref{prop:pdf-timelike-dglap}、\cref{prop:coll-photon-sf}、\cref{prop:coll-luxqed}、\cref{thm:coll-g1-moments}、\cref{prop:coll-transversity-evolution}、\cref{thm:tmd-sivers-sign}、\cref{prop:excl-erbl}、\cref{prop:excl-polynomiality}、\cref{thm:heavy-inclusive-b}、\cref{thm:lamet-exp-decay}、\cref{lem:appM-forward-tube}、\cref{lem:appM-positivity}。
解除途径：测试解除途径。
% uses: none
\end{assumption}

\begin{assumption}[规范不变复合算符与光线算符的存在与重整化群 \tagC]\label{asm:intro-composite}
（测试用合成陈述：内容从略。）
非微扰地位：\cref{asm:intro-wightman}、\cref{thm:prelim-composite-renorm}、\cref{thm:pdf-moments-local}、\cref{asm:intro-np-ope}、\cref{prop:dis-ope-rg}、\cref{thm:dis-moments-np}、\cref{cor:dis-weak-xspace}、\cref{thm:dis-moments}、\cref{prop:dis-twist4}、\cref{prop:dis-tmc}、\cref{thm:dis-bjorken}、\cref{prop:pdf-invariance}、\cref{prop:pdf-support-hadron}、\cref{thm:pdf-dglap}、\cref{thm:pdf-dis-xspace}、\cref{prop:pdf-sum-rules}、\cref{prop:pdf-timelike-dglap}、\cref{prop:coll-photon-sf}、\cref{prop:coll-luxqed}、\cref{thm:coll-g1-moments}、\cref{prop:coll-transversity-evolution}、\cref{prop:tmd-small-b}、\cref{thm:tmd-sivers-sign}、\cref{prop:excl-erbl}、\cref{prop:excl-polynomiality}、\cref{thm:heavy-inclusive-b}、\cref{lem:appM-norm-bound}、\cref{def:intro-MA}。
地位：测试状态句。
文献：测试。
被使用于：\cref{asm:intro-np-ope}、\cref{def:intro-MA}、\cref{prop:dis-ope-rg}、\cref{thm:dis-moments-np}、\cref{cor:dis-weak-xspace}、\cref{thm:dis-moments}、\cref{prop:dis-twist4}、\cref{prop:dis-tmc}、\cref{thm:dis-bjorken}、\cref{prop:pdf-invariance}、\cref{prop:pdf-support-hadron}、\cref{thm:pdf-moments-local}、\cref{thm:pdf-dglap}、\cref{thm:pdf-dis-xspace}、\cref{prop:pdf-sum-rules}、\cref{prop:pdf-timelike-dglap}、\cref{prop:coll-photon-sf}、\cref{prop:coll-luxqed}、\cref{thm:coll-g1-moments}、\cref{prop:coll-transversity-evolution}、\cref{prop:tmd-small-b}、\cref{thm:tmd-sivers-sign}、\cref{prop:excl-erbl}、\cref{prop:excl-polynomiality}、\cref{thm:heavy-inclusive-b}、\cref{lem:appM-norm-bound}。
解除途径：测试解除途径。
\uses{asm:intro-wightman}
\end{assumption}

\section{第 6 节（测试）}\label{sec:intro-overview}

\begin{assumption}[非微扰算符乘积展开 \tagC]\label{asm:intro-np-ope}
（测试用合成陈述：内容从略。）
非微扰地位：\cref{conj:ee-np-ope}、\cref{asm:intro-af-coeff}、\cref{lem:dis-ope-segment}、\cref{thm:dis-fixed-n}、\cref{thm:dis-moments-np}、\cref{cor:dis-weak-xspace}、\cref{thm:dis-moments}、\cref{prop:dis-twist4}、\cref{prop:dis-tmc}、\cref{thm:dis-bjorken}、\cref{thm:pdf-dis-xspace}、\cref{prop:coll-photon-sf}、\cref{thm:coll-g1-moments}、\cref{def:intro-MA}。
地位：测试状态句。
文献：测试。
被使用于：\cref{asm:intro-af-coeff}、\cref{def:intro-MA}、\cref{lem:dis-ope-segment}、\cref{thm:dis-fixed-n}、\cref{thm:dis-moments-np}、\cref{cor:dis-weak-xspace}、\cref{thm:dis-moments}、\cref{prop:dis-twist4}、\cref{prop:dis-tmc}、\cref{thm:dis-bjorken}、\cref{thm:pdf-dis-xspace}、\cref{prop:coll-photon-sf}、\cref{thm:coll-g1-moments}。
解除途径：测试解除途径。
\uses{asm:intro-wightman, asm:intro-composite}
\end{assumption}

\begin{assumption}[系数的渐近自由行为 \tagC]\label{asm:intro-af-coeff}
（测试用合成陈述：内容从略。）
非微扰地位：\cref{thm:dis-moments-np}、\cref{cor:dis-weak-xspace}、\cref{thm:dis-moments}、\cref{prop:dis-twist4}、\cref{prop:dis-tmc}、\cref{thm:dis-bjorken}、\cref{thm:pdf-dis-xspace}、\cref{prop:coll-photon-sf}、\cref{thm:coll-g1-moments}、\cref{thm:heavy-inclusive-b}、\cref{lem:appM-norm-bound}、\cref{def:intro-MA}、\cref{asm:intro-np-ope}。
地位：测试状态句。
文献：测试。
被使用于：\cref{def:intro-MA}、\cref{thm:dis-moments-np}、\cref{cor:dis-weak-xspace}、\cref{thm:dis-moments}、\cref{prop:dis-twist4}、\cref{prop:dis-tmc}、\cref{thm:dis-bjorken}、\cref{thm:pdf-dis-xspace}、\cref{prop:coll-photon-sf}、\cref{thm:coll-g1-moments}、\cref{thm:heavy-inclusive-b}、\cref{lem:appM-norm-bound}。
解除途径：测试解除途径。
\uses{asm:intro-np-ope}
\end{assumption}

\section{第 7 节（测试）}\label{sec:intro-misconceptions}

\begin{assumption}[前向振幅的增长条件 \tagC]\label{asm:intro-growth}
（测试用合成陈述：内容从略。）
非微扰地位：\cref{prop:dis-low-moments}、\cref{conj:coll-bc}、\cref{thm:dis-moments-np}、\cref{thm:dis-moments}、\cref{thm:dis-bjorken}、\cref{thm:pdf-dis-xspace}、\cref{def:intro-MA}。
地位：测试状态句。
文献：测试。
被使用于：\cref{def:intro-MA}、\cref{thm:dis-moments-np}、\cref{thm:dis-moments}、\cref{thm:dis-bjorken}、\cref{thm:pdf-dis-xspace}。
解除途径：测试解除途径。
\uses{asm:intro-wightman}
\end{assumption}

\begin{assumption}[从部分子微扰论到强子的桥接假设 \tagC]\label{asm:intro-bridge}
（测试用合成陈述：内容从略。）
非微扰地位：\cref{thm:method-template}、\cref{thm:disall-factorization}、\cref{thm:disall-heavy}、\cref{prop:disall-nuclear}、\cref{thm:coll-sia}、\cref{thm:coll-sidis}、\cref{thm:coll-direct-photo}、\cref{thm:dy-factorization}、\cref{thm:dy-colour-singlet}、\cref{thm:dy-polarized}、\cref{thm:dy-higher-twist}、\cref{thm:tmd-factorization}、\cref{prop:tmd-small-b}、\cref{cor:tmd-sivers-observable}、\cref{prop:tmd-collins-universality}、\cref{thm:scet-thrust}、\cref{thm:jets-ngl-factorization}、\cref{thm:break-ddis}、\cref{thm:break-double-dy}、\cref{thm:excl-pion-ff}、\cref{thm:excl-transition}、\cref{thm:excl-dvcs}、\cref{thm:excl-meson-L}、\cref{thm:heavy-bdpi}、\cref{prop:heavy-btopi-ff}、\cref{prop:heavy-qcdf-lp}、\cref{thm:heavy-inclusive-b}、\cref{thm:heavy-nrqcd-decay}、\cref{prop:heavy-quarkonium-fragmentation}、\cref{prop:power-dis-twist4-x}、\cref{def:intro-MA}。
地位：测试状态句。
文献：测试。
被使用于：\cref{def:intro-MA}、\cref{thm:method-template}、\cref{thm:disall-factorization}、\cref{thm:disall-heavy}、\cref{prop:disall-nuclear}、\cref{thm:coll-sia}、\cref{thm:coll-sidis}、\cref{thm:coll-direct-photo}、\cref{thm:dy-factorization}、\cref{thm:dy-colour-singlet}、\cref{thm:dy-polarized}、\cref{thm:dy-higher-twist}、\cref{thm:tmd-factorization}、\cref{prop:tmd-small-b}、\cref{cor:tmd-sivers-observable}、\cref{prop:tmd-collins-universality}、\cref{thm:scet-thrust}、\cref{thm:jets-ngl-factorization}、\cref{thm:break-ddis}、\cref{thm:break-double-dy}、\cref{thm:excl-pion-ff}、\cref{thm:excl-transition}、\cref{thm:excl-dvcs}、\cref{thm:excl-meson-L}、\cref{thm:heavy-bdpi}、\cref{prop:heavy-btopi-ff}、\cref{prop:heavy-qcdf-lp}、\cref{thm:heavy-inclusive-b}、\cref{thm:heavy-nrqcd-decay}、\cref{prop:heavy-quarkonium-fragmentation}、\cref{prop:power-dis-twist4-x}。
解除途径：测试解除途径。
\uses{def:intro-factorization}
\end{assumption}

\section{第 8 节（测试）}\label{sec:intro-organisation}

\begin{definition}[元假设 MA0–MA4（总称）与编号假设的对应]\label{def:intro-MA}
（测试用合成陈述：内容从略。）
\uses{asm:intro-wightman, asm:intro-composite, asm:intro-np-ope, asm:intro-af-coeff, asm:intro-growth, asm:intro-bridge}
\end{definition}
